% Options for packages loaded elsewhere
\PassOptionsToPackage{unicode}{hyperref}
\PassOptionsToPackage{hyphens}{url}
\documentclass[
  11pt,
  a4paper,
]{article}
\usepackage{xcolor}
\usepackage[margin=18mm]{geometry}
\usepackage{amsmath,amssymb}
\setcounter{secnumdepth}{-\maxdimen} % remove section numbering
\usepackage{iftex}
\ifPDFTeX
  \usepackage[T1]{fontenc}
  \usepackage[utf8]{inputenc}
  \usepackage{textcomp} % provide euro and other symbols
\else % if luatex or xetex
  \usepackage{unicode-math} % this also loads fontspec
  \defaultfontfeatures{Scale=MatchLowercase}
  \defaultfontfeatures[\rmfamily]{Ligatures=TeX,Scale=1}
\fi
\ifPDFTeX\else
  % xetex/luatex font selection
    \setmainfont[]{Noto Serif}
    \setsansfont[]{Noto Sans}
    \setmonofont[]{Noto Sans Mono}
  \setmathfont[]{Noto Sans Math}
\fi
% Use upquote if available, for straight quotes in verbatim environments
\IfFileExists{upquote.sty}{\usepackage{upquote}}{}
\IfFileExists{microtype.sty}{% use microtype if available
  \usepackage[]{microtype}
  \UseMicrotypeSet[protrusion]{basicmath} % disable protrusion for tt fonts
}{}
\makeatletter
\@ifundefined{KOMAClassName}{% if non-KOMA class
  \IfFileExists{parskip.sty}{%
    \usepackage{parskip}
  }{% else
    \setlength{\parindent}{0pt}
    \setlength{\parskip}{6pt plus 2pt minus 1pt}}
}{% if KOMA class
  \KOMAoptions{parskip=half}}
\makeatother
\ifLuaTeX
\usepackage[bidi=basic]{babel}
\else
\usepackage[bidi=default]{babel}
\fi
\babelprovide[main,import]{russian}
\ifPDFTeX
\else
\babelfont{rm}[]{Noto Serif}
\fi
% get rid of language-specific shorthands (see #6817):
\let\LanguageShortHands\languageshorthands
\def\languageshorthands#1{}
\setlength{\emergencystretch}{3em} % prevent overfull lines
\providecommand{\tightlist}{%
  \setlength{\itemsep}{0pt}\setlength{\parskip}{0pt}}
\usepackage{bookmark}
\IfFileExists{xurl.sty}{\usepackage{xurl}}{} % add URL line breaks if available
\urlstyle{same}
\hypersetup{
  pdftitle={Билеты по дискретным случайным процессам},
  pdflang={ru-RU},
  hidelinks,
  pdfcreator={LaTeX via pandoc}}

\title{Билеты по дискретным случайным процессам}
\author{}
\date{}

\begin{document}
\maketitle

\section{Билет 25. Спектральные
характеристики}\label{ux431ux438ux43bux435ux442-25.-ux441ux43fux435ux43aux442ux440ux430ux43bux44cux43dux44bux435-ux445ux430ux440ux430ux43aux442ux435ux440ux438ux441ux442ux438ux43aux438}

Источники: Ширяев 2, гл. VI (Спектральная теория); лекции ДСП;
Сердобольская.

\subsection{1. Короткий
ответ}\label{ux43aux43eux440ux43eux442ux43aux438ux439-ux43eux442ux432ux435ux442}

Любая стационарная в широком смысле случайная последовательность (с
дискретным временем)

\[
X = (X_n)_{n \in \mathbb{Z}}
\]

и её ковариационная функция

\[
R(k) = \operatorname{Cov}(X_{n+k}, X_n)
\]

допускают спектральное представление. Это значит, что процесс можно
представить как суперпозицию некоррелированных гармонических колебаний с
различными частотами

\[
\lambda \in [-\pi, \pi].
\]

Главный результат для ковариационной функции --- \textbf{теорема
Герглотца}: любая ковариационная функция \(R(k)\) представляется как
преобразование Фурье некоторой конечной меры на \([-\pi, \pi]\)
(спектральной меры \(F\)):

\[
R(k) = \int_{-\pi}^{\pi} e^{i k \lambda} F(d\lambda).
\]

Если мера \(F\) имеет плотность \(f(\lambda)\) по мере Лебега, то она
называется \textbf{спектральной плотностью}, и \(R(k)\) связана с ней
взаимно обратными преобразованиями Фурье.

Сам процесс \(X_n\) также представляется интегралом:

\[
X_n = \int_{-\pi}^{\pi} e^{i n \lambda} Z(d\lambda),
\]

где \(Z\) --- случайная мера с ортогональными приращениями. Это
представление изометрично связывает гильбертово пространство величин
процесса \(H_X\) (см. Билет 12) и пространство функций \(L_2(F)\).

\subsection{2. Полный
ответ}\label{ux43fux43eux43bux43dux44bux439-ux43eux442ux432ux435ux442}

\subsubsection{2.1. Введение: ковариационная функция и её
свойства}\label{ux432ux432ux435ux434ux435ux43dux438ux435-ux43aux43eux432ux430ux440ux438ux430ux446ux438ux43eux43dux43dux430ux44f-ux444ux443ux43dux43aux446ux438ux44f-ux438-ux435ux451-ux441ux432ux43eux439ux441ux442ux432ux430}

Рассмотрим комплекснозначную случайную последовательность со вторыми
моментами

\[
\mathbb E|X_n|^2 < \infty.
\]

Как было показано в Билете 17, процесс стационарен в широком смысле,
если его математическое ожидание постоянно:

\[
\mathbb E X_n = m,
\]

а ковариация зависит только от сдвига:

\[
\operatorname{Cov}(X_{n+k}, X_n) = R(k).
\]

Без ограничения общности далее будем считать, что процесс центрирован:

\[
m = 0, \quad R(k) = \mathbb E[X_k \overline{X_0}].
\]

В Билете 24 было доказано, что ковариационная функция неотрицательно
определена: для любых комплексных чисел \(a_1, \ldots, a_r\) и моментов
времени \(t_1, \ldots, t_r\) выполнено

\[
\sum_{j=1}^r \sum_{l=1}^r a_j \overline{a_l} R(t_j - t_l) = \mathbb E \left| \sum_{j=1}^r a_j X_{t_j} \right|^2 \ge 0.
\]

Это свойство роднит ковариационную функцию с характеристическими
функциями распределений (см. Билет 08). Для последних работает теорема
Бохнера---Хинчина, утверждающая, что непрерывная положительно
определенная функция является преобразованием Фурье вероятностной меры.
Аналогичный результат для дискретного времени называется теоремой
Герглотца.

\subsubsection{2.2. Теорема Герглотца о
спектре}\label{ux442ux435ux43eux440ux435ux43cux430-ux433ux435ux440ux433ux43bux43eux442ux446ux430-ux43e-ux441ux43fux435ux43aux442ux440ux435}

Так как время \(t \in \mathbb{Z}\) дискретно, экспоненты
\(e^{i t \lambda}\) периодичны по частоте \(\lambda\) с периодом
\(2\pi\). Поэтому ``весь спектр'' частот, доступный для процесса с
дискретным временем, помещается на одном периоде, например на
полуинтервале

\[
[-\pi, \pi).
\]

\textbf{Теорема Герглотца} утверждает, что \(R(k)\) неотрицательно
определена тогда и только тогда, когда она является преобразованием
Фурье-Стилтьеса:

\[
R(k) = \int_{-\pi}^{\pi} e^{i k \lambda} F(d\lambda),
\]

где \(F\) --- конечная мера на борелевских множествах \([-\pi, \pi)\),
называемая \textbf{спектральной мерой}. Выбор \([-\pi,\pi)\) вместо
\([-\pi,\pi]\) нужен только для того, чтобы не считать одну и ту же
частоту дважды: точки \(-\pi\) и \(\pi\) дают одинаковые значения
\(e^{ik\lambda}\) при целых \(k\).

Физический смысл спектральной меры: дисперсия процесса (полная мощность)
равна

\[
\operatorname{Var}(X_n) = R(0) = \int_{-\pi}^{\pi} F(d\lambda) = F([-\pi, \pi)).
\]

Мера интервала \(F([\lambda_1, \lambda_2])\) показывает вклад этих
частот в полную мощность процесса. Доля мощности равна
\(F([\lambda_1,\lambda_2])/R(0)\) при \(R(0)>0\).

Часто вместо меры пишут \textbf{спектральную функцию}

\[
\mathcal F(\lambda)=F([-\pi,\lambda]), \qquad -\pi\le \lambda<\pi.
\]

Это не новая сущность, а способ записать ту же меру через неубывающую
функцию. Тогда интеграл выше читается как интеграл Лебега-Стилтьеса:

\[
R(k)=\int_{-\pi}^{\pi}e^{ik\lambda}\,d\mathcal F(\lambda).
\]

Спектральная мера единственна: если две конечные меры имеют одинаковые
коэффициенты Фурье \(\int e^{ik\lambda}F(d\lambda)\) для всех
\(k\in\mathbb Z\), то они совпадают. Идея доказательства единственности:
тригонометрические полиномы равномерно приближают непрерывные функции на
окружности, значит совпадают интегралы от всех непрерывных функций, а
значит совпадают и меры.

\subsubsection{2.3. Спектральная
плотность}\label{ux441ux43fux435ux43aux442ux440ux430ux43bux44cux43dux430ux44f-ux43fux43bux43eux442ux43dux43eux441ux442ux44c}

В общем случае по теореме Лебега о разложении (связано с Билетом 10),
спектральная мера распадается на три компоненты: абсолютно непрерывную,
дискретную (атомарную) и сингулярную:

\[
F = F_{ac} + F_d + F_{sing}.
\]

Если мера абсолютно непрерывна относительно меры Лебега, существует
функция \(f(\lambda) \ge 0\), называемая \textbf{спектральной
плотностью}, такая что:

\[
F(d\lambda) = f(\lambda) d\lambda.
\]

В этом случае теорема Герглотца принимает вид:

\[
R(k) = \int_{-\pi}^{\pi} e^{i k \lambda} f(\lambda) d\lambda.
\]

Если ряд \(\sum_{k=-\infty}^{\infty} |R(k)|\) сходится, то спектральная
плотность существует, непрерывна, и её можно найти по формуле обращения:

\[
f(\lambda) = \frac{1}{2\pi} \sum_{k=-\infty}^{\infty} R(k) e^{-i k \lambda}.
\]

В более слабой форме, если \(\sum_k |R(k)|^2<\infty\), то спектральная
функция тоже имеет плотность, но уже в смысле \(L_2\):

\[
f(\lambda)=\frac1{2\pi}\sum_{k=-\infty}^{\infty}R(k)e^{-ik\lambda}
\quad \text{в } L_2[-\pi,\pi].
\]

Для экзамена удобно помнить иерархию: абсолютная суммируемость
ковариаций дает хорошую непрерывную плотность; квадрат-суммируемость
дает плотность как \(L_2\)-функцию; без этих условий остается общая
спектральная мера, у которой могут быть атомы и сингулярная часть.

Для вещественного процесса \(R(k)\) вещественна и четна
(\(R(k) = R(-k)\)). Поэтому спектральная плотность также является четной
функцией:

\[
f(\lambda) = f(-\lambda),
\]

а синусы в формуле Эйлера взаимно уничтожаются:

\[
f(\lambda) = \frac{1}{2\pi} R(0) + \frac{1}{\pi} \sum_{k=1}^{\infty} R(k) \cos(k \lambda).
\]

\subsubsection{2.4. Спектральное представление стационарной
последовательности}\label{ux441ux43fux435ux43aux442ux440ux430ux43bux44cux43dux43eux435-ux43fux440ux435ux434ux441ux442ux430ux432ux43bux435ux43dux438ux435-ux441ux442ux430ux446ux438ux43eux43dux430ux440ux43dux43eux439-ux43fux43eux441ux43bux435ux434ux43eux432ux430ux442ux435ux43bux44cux43dux43eux441ux442ux438}

Мы выразили через частоты функцию ковариации. Можно ли выразить через
гармоники сам случайный процесс?

Да. Любую центрированную стационарную в широком смысле
последовательность можно записать в виде \textbf{стохастического
интеграла}:

\[
X_n = \int_{-\pi}^{\pi} e^{i n \lambda} Z(d\lambda),
\]

где \(Z\) --- случайная мера с ортогональными приращениями.

Что значит ``случайная мера с ортогональными приращениями''? Каждому
борелевскому множеству \(\Delta \subset [-\pi, \pi)\) ставится в
соответствие комплекснозначная случайная величина
\(Z(\Delta) \in L_2(\Omega)\). 1. Она центрирована:
\(\mathbb E Z(\Delta) = 0\). 2. Ее дисперсия задается спектральной мерой
\(F\):

\[
\mathbb E|Z(\Delta)|^2 = F(\Delta).
\]

\begin{enumerate}
\def\labelenumi{\arabic{enumi}.}
\setcounter{enumi}{2}
\tightlist
\item
  Для непересекающихся \(\Delta_1\) и \(\Delta_2\) величины
  \(Z(\Delta_1)\) и \(Z(\Delta_2)\) некоррелированы (ортогональны в
  пространстве \(L_2(\Omega)\)):
\end{enumerate}

\[
\mathbb E[Z(\Delta_1) \overline{Z(\Delta_2)}] = 0.
\]

\subsubsection{2.5. Изометрия гильбертовых
пространств}\label{ux438ux437ux43eux43cux435ux442ux440ux438ux44f-ux433ux438ux43bux44cux431ux435ux440ux442ux43eux432ux44bux445-ux43fux440ux43eux441ux442ux440ux430ux43dux441ux442ux432}

Понятие ортогональной случайной меры и интеграла по ней напрямую
опирается на теорию гильбертовых пространств (см. Билет 12).

Пусть \(H_X\) --- замкнутое линейное подпространство в \(L_2(\Omega)\),
порожденное величинами \(\{X_n\}_{n \in \mathbb{Z}}\). А \(L_2(F)\) ---
пространство детерминированных (не случайных!) комплекснозначных функций
на \([-\pi, \pi]\), интегрируемых с квадратом по мере \(F\).

Интеграл \(\int g(\lambda) Z(d\lambda)\) строится как предел линейных
комбинаций: мы берем ступенчатую функцию
\(g(\lambda) = \sum c_j \mathbf{1}_{\Delta_j}(\lambda)\) и определяем
для нее интеграл как \(\sum c_j Z(\Delta_j)\).

Из свойства ортогональности приращений \(Z\) следует фундаментальное
свойство: скалярное произведение случайных величин (в \(H_X\)) в
точности равно скалярному произведению функций (в \(L_2(F)\)).

Действительно, для ступенчатых функций:

\[
\mathbb E \left[ \left( \sum_j c_j Z(\Delta_j) \right) \overline{\left( \sum_k d_k Z(\Delta_k) \right)} \right]
=
\sum_j c_j \overline{d_j} F(\Delta_j)
=
\int_{-\pi}^{\pi} g(\lambda) \overline{h(\lambda)} F(d\lambda).
\]

Таким образом, отображение

\[
H_X \ni \int_{-\pi}^{\pi} g(\lambda) Z(d\lambda) \quad \longleftrightarrow \quad g(\lambda) \in L_2(F)
\]

является \textbf{изометрическим изоморфизмом} между гильбертовым
пространством случайных величин и гильбертовым пространством функций.

Элементу \(X_n \in H_X\) соответствует функция
\(e^{i n \lambda} \in L_2(F)\). Скалярное произведение
\(\langle X_n, X_0 \rangle_{H_X} = R(n)\) переходит в:

\[
\int_{-\pi}^{\pi} e^{i n \lambda} \cdot 1 \, F(d\lambda),
\]

что в точности возвращает нас к теореме Герглотца.

\subsubsection{2.6. Связь спектра с усреднением и
эргодичностью}\label{ux441ux432ux44fux437ux44c-ux441ux43fux435ux43aux442ux440ux430-ux441-ux443ux441ux440ux435ux434ux43dux435ux43dux438ux435ux43c-ux438-ux44dux440ux433ux43eux434ux438ux447ux43dux43eux441ux442ux44cux44e}

Спектр сразу объясняет результат из Билета 24 про временные средние. Для
центрированного процесса

\[
\overline X_N=\frac1N\sum_{n=0}^{N-1}X_n
\]

по спектральному представлению:

\[
\overline X_N
=
\int_{-\pi}^{\pi}
\left(\frac1N\sum_{n=0}^{N-1}e^{in\lambda}\right)
Z(d\lambda).
\]

По изометрии:

\[
\mathbb E|\overline X_N|^2
=
\int_{-\pi}^{\pi}
\left|
\frac1N\sum_{n=0}^{N-1}e^{in\lambda}
\right|^2
F(d\lambda).
\]

Множитель под интегралом является нормированным ядром Фейера: при
\(\lambda\ne0\) он стремится к нулю, а в точке \(\lambda=0\) равен
\(1\). Поэтому предел дисперсии среднего равен массе спектральной меры в
нуле:

\[
\mathbb E|\overline X_N|^2 \to F(\{0\}).
\]

Значит, эргодичность по среднему в широком смысле эквивалентна
отсутствию атома в нулевой частоте:

\[
F(\{0\})=0.
\]

Интуитивно атом в нуле --- это случайная постоянная компонента процесса.
Она не исчезает при усреднении по времени.

\subsection{3. Основные
определения}\label{ux43eux441ux43dux43eux432ux43dux44bux435-ux43eux43fux440ux435ux434ux435ux43bux435ux43dux438ux44f}

\textbf{Спектральная мера} \(F\) --- конечная мера на борелевской
\(\sigma\)-алгебре одного периода частот, обычно \([-\pi, \pi)\), такая
что ковариационная функция \(R(k)\) является ее преобразованием
Фурье-Стилтьеса.

\textbf{Спектральная функция} \(\mathcal F(\lambda)\) --- функция
распределения спектральной меры:

\[
\mathcal F(\lambda)=F([-\pi,\lambda]).
\]

Она неубывающая и правая непрерывная; скачок
\(\mathcal F(\lambda)-\mathcal F(\lambda-)\) равен массе атома в частоте
\(\lambda\).

\textbf{Спектральная плотность} \(f(\lambda)\) --- производная
Радона-Никодима (см. Билет 10) спектральной меры относительно меры
Лебега, существующая, если мера \(F\) абсолютно непрерывна:

\[
F(d\lambda) = f(\lambda) d\lambda.
\]

\textbf{Случайная мера с ортогональными приращениями} \(Z\) --- функция
множества, заданная на борелевской \(\sigma\)-алгебре \([-\pi, \pi)\),
значениями которой являются случайные величины из \(L_2(\Omega)\),
удовлетворяющая: 1. \(\mathbb E Z(\Delta) = 0\). 2.
\(\Delta_1 \cap \Delta_2 = \emptyset \implies Z(\Delta_1 \cup \Delta_2) = Z(\Delta_1) + Z(\Delta_2)\)
почти наверное. 3.
\(\Delta_1 \cap \Delta_2 = \emptyset \implies \mathbb E[Z(\Delta_1) \overline{Z(\Delta_2)}] = 0\).

\textbf{Структурная мера} \(\mu_Z\) для меры \(Z\) --- детерминированная
мера \(\mu_Z(\Delta) = \mathbb E|Z(\Delta)|^2\). (В контексте
стационарных процессов структурная мера совпадает со спектральной мерой
\(F\)).

\textbf{Частотная характеристика фильтра} с импульсной характеристикой
\(h_k\) --- функция

\[
H(\lambda)=\sum_{k=-\infty}^{\infty}h_k e^{-ik\lambda},
\]

которая показывает, как линейный фильтр усиливает или подавляет разные
частоты.

\subsection{4. Основные теоремы и
свойства}\label{ux43eux441ux43dux43eux432ux43dux44bux435-ux442ux435ux43eux440ux435ux43cux44b-ux438-ux441ux432ux43eux439ux441ux442ux432ux430}

\subsubsection{Теорема
Герглотца}\label{ux442ux435ux43eux440ux435ux43cux430-ux433ux435ux440ux433ux43bux43eux442ux446ux430}

\textbf{Условие:} Последовательность чисел \(R(k)\),
\(k \in \mathbb{Z}\), является ковариационной функцией стационарной в
широком смысле последовательности. \textbf{Утверждение:} Это выполнено
тогда и только тогда, когда существует конечная мера \(F\) на
\([-\pi, \pi)\), такая что:

\[
R(k) = \int_{-\pi}^{\pi} e^{i k \lambda} F(d\lambda).
\]

Мера \(F\) определяется функцией \(R\) единственным образом. Кроме того,

\[
F([-\pi,\pi))=R(0),
\]

то есть полная масса спектральной меры равна дисперсии центрированного
процесса.

\textbf{Схема доказательства:} 1. Необходимость. Так как \(R(k)\)
положительно определена на решетке \(\mathbb{Z}\), строят сглаженные
неотрицательные тригонометрические многочлены:

\[
f_N(\lambda) = \frac{1}{2\pi N} \sum_{n=0}^{N-1} \sum_{m=0}^{N-1} R(n-m) e^{-i(n-m)\lambda} \ge 0.
\]

Положительность следует из неотрицательной определенности \(R\).
Интеграл от \(f_N(\lambda)\) равен \(R(0)\). Меры
\(F_N(d\lambda) = f_N(\lambda) d\lambda\) образуют слабо компактное
семейство (по теореме Прохорова). Их слабый предел \(F\) и будет искомой
мерой.

\begin{enumerate}
\def\labelenumi{\arabic{enumi}.}
\setcounter{enumi}{1}
\tightlist
\item
  Достаточность. Если интеграл существует, то положительная
  определенность \(R(k)\) очевидна:
\end{enumerate}

\[
\sum_{n,m} a_n \overline{a_m} R(n-m) = \int_{-\pi}^{\pi} \left| \sum_n a_n e^{i n \lambda} \right|^2 F(d\lambda) \ge 0.
\]

По теореме Колмогорова (см. Билет 15) существование неотрицательно
определенной матрицы всегда гарантирует существование гауссовского
стационарного процесса с такой ковариацией.

\subsubsection{Теорема Карунена о спектральном
представлении}\label{ux442ux435ux43eux440ux435ux43cux430-ux43aux430ux440ux443ux43dux435ux43dux430-ux43e-ux441ux43fux435ux43aux442ux440ux430ux43bux44cux43dux43eux43c-ux43fux440ux435ux434ux441ux442ux430ux432ux43bux435ux43dux438ux438}

\textbf{Условие:} Пусть \(X_n\) --- центрированная стационарная в
широком смысле последовательность со спектральной мерой \(F\).
\textbf{Утверждение:} Существует случайная мера \(Z\) с ортогональными
приращениями, такая что \(\mathbb E|Z(\Delta)|^2 = F(\Delta)\) и

\[
X_n = \int_{-\pi}^{\pi} e^{i n \lambda} Z(d\lambda).
\]

\textbf{Схема доказательства:} Зададим отображение
\(I: e^{i n \lambda} \mapsto X_n\). Оно сохраняет скалярное
произведение, так как
\(\langle e^{i n \lambda}, e^{i m \lambda} \rangle_{L_2(F)} = \int e^{i(n-m)\lambda}F(d\lambda) = R(n-m) = \langle X_n, X_m \rangle_{H_X}\).
По линейности продолжим \(I\) на тригонометрические полиномы. Множество
тригонометрических полиномов плотно в \(L_2(F)\). Значит, мы можем по
непрерывности продолжить \(I\) на все пространство \(L_2(F)\). Определим
меру \(Z\) как образы индикаторов:

\[
Z(\Delta) = I(\mathbf{1}_\Delta(\lambda)).
\]

Тогда \(Z(\Delta)\) ортогональны для непересекающихся множеств (т.к.
ортогональны их индикаторы в \(L_2(F)\)), а \(X_n\) получается
интегрированием функции \(e^{in\lambda}\) по этой мере.

\subsubsection{Влияние линейных
фильтров}\label{ux432ux43bux438ux44fux43dux438ux435-ux43bux438ux43dux435ux439ux43dux44bux445-ux444ux438ux43bux44cux442ux440ux43eux432}

Если пропустить процесс \(X_n\) через устойчивый линейный фильтр с
импульсной характеристикой \(h_k\) (где \(\sum |h_k| < \infty\)):

\[
Y_n = \sum_{k=-\infty}^{\infty} h_k X_{n-k},
\]

то спектральная плотность преобразуется путем умножения на квадрат
модуля передаточной функции (частотной характеристики) фильтра:

\[
f_Y(\lambda) = |H(\lambda)|^2 f_X(\lambda), \quad \text{где } H(\lambda) = \sum_{k=-\infty}^{\infty} h_k e^{-i k \lambda}.
\]

\subsection{5. Примеры}\label{ux43fux440ux438ux43cux435ux440ux44b}

\subsubsection{Пример 1. Белый шум (процесс без
памяти)}\label{ux43fux440ux438ux43cux435ux440-1.-ux431ux435ux43bux44bux439-ux448ux443ux43c-ux43fux440ux43eux446ux435ux441ux441-ux431ux435ux437-ux43fux430ux43cux44fux442ux438}

Пусть \(X_n = \varepsilon_n\), где \(\varepsilon_n\) ---
некоррелированные величины:

\[
\mathbb E \varepsilon_n = 0, \quad \operatorname{Var}(\varepsilon_n) = \sigma^2, \quad \operatorname{Cov}(\varepsilon_n, \varepsilon_m) = 0 \text{ при } n \ne m.
\]

Ковариационная функция:

\[
R(0) = \sigma^2, \quad R(k) = 0 \text{ для } k \ne 0.
\]

Ряд Фурье для плотности:

\[
f(\lambda) = \frac{1}{2\pi} \sum_{k=-\infty}^{\infty} R(k) e^{-i k \lambda} = \frac{\sigma^2}{2\pi}.
\]

\textbf{Вывод:} Спектральная плотность белого шума константна. В
процессе присутствуют все частоты от \(-\pi\) до \(\pi\) в равной
пропорции. Отсюда аналогия с белым светом в оптике.

\subsubsection{Пример 2. Гармоническое колебание со случайной фазой
(линейчатый
спектр)}\label{ux43fux440ux438ux43cux435ux440-2.-ux433ux430ux440ux43cux43eux43dux438ux447ux435ux441ux43aux43eux435-ux43aux43eux43bux435ux431ux430ux43dux438ux435-ux441ux43e-ux441ux43bux443ux447ux430ux439ux43dux43eux439-ux444ux430ux437ux43eux439-ux43bux438ux43dux435ux439ux447ux430ux442ux44bux439-ux441ux43fux435ux43aux442ux440}

Пусть

\[
X_n = A \cos(\omega n + \Phi),
\]

где \(\omega \in (0, \pi)\) и \(A > 0\) --- детерминированные числа, а
\(\Phi \sim U[0, 2\pi]\) --- случайная величина.

Считаем ковариацию:

\[
R(k) = \mathbb E[A \cos(\omega(n+k) + \Phi) A \cos(\omega n + \Phi)].
\]

Используя формулу произведения косинусов, получаем:

\[
R(k) = \frac{A^2}{2} \cos(\omega k).
\]

Представим косинус через экспоненты:

\[
R(k) = \frac{A^2}{4} e^{i \omega k} + \frac{A^2}{4} e^{-i \omega k}.
\]

\textbf{Вывод:} Это уже преобразование Фурье для меры \(F\), которая
состоит ровно из двух атомов (дельта-функций Дирака, см. Билет 20) в
точках \(\lambda_1 = \omega\) и \(\lambda_2 = -\omega\), масса каждого
из которых равна \(\frac{A^2}{4}\). Плотности \(f(\lambda)\) по мере
Лебега у такого процесса \textbf{не существует}. Весь ``спектр''
сосредоточен в двух конкретных частотах.

\subsubsection{Пример 3. Авторегрессия первого порядка
AR(1)}\label{ux43fux440ux438ux43cux435ux440-3.-ux430ux432ux442ux43eux440ux435ux433ux440ux435ux441ux441ux438ux44f-ux43fux435ux440ux432ux43eux433ux43e-ux43fux43eux440ux44fux434ux43aux430-ar1}

Пусть процесс задан разностным уравнением:

\[
X_n = a X_{n-1} + \varepsilon_n, \quad |a| < 1,
\]

где \(\varepsilon_n\) --- белый шум с дисперсией \(\sigma^2\).

Как мы знаем, его ковариационная функция равна:

\[
R(k) = \frac{\sigma^2}{1-a^2} a^{|k|}.
\]

Найдем спектральную плотность. Так как процесс получается из белого шума
фильтром \(H(\lambda) = \frac{1}{1 - a e^{-i\lambda}}\) (это следует из
\(X_n - a X_{n-1} = \varepsilon_n\), т.е.
\(X(1 - a e^{-i\lambda}) = \varepsilon\)), то по свойству фильтра:

\[
f_X(\lambda) = |H(\lambda)|^2 f_\varepsilon(\lambda) = \left| \frac{1}{1 - a e^{-i\lambda}} \right|^2 \frac{\sigma^2}{2\pi}.
\]

Распишем знаменатель:

\[
|1 - a \cos\lambda + i a \sin\lambda|^2 = (1 - a \cos\lambda)^2 + (a \sin\lambda)^2 = 1 - 2a \cos\lambda + a^2.
\]

\textbf{Вывод:}

\[
f(\lambda) = \frac{\sigma^2}{2\pi (1 - 2a \cos\lambda + a^2)}.
\]

Если \(a > 0\), спектр имеет пик на низких частотах (возле
\(\lambda = 0\)), процесс имеет длинные ``волны''. Если \(a < 0\), пик
на высоких частотах (возле \(\lambda = \pi\)), процесс ``мелко дрожит''.

\subsubsection{Пример 4. Скользящее среднее
MA(1)}\label{ux43fux440ux438ux43cux435ux440-4.-ux441ux43aux43eux43bux44cux437ux44fux449ux435ux435-ux441ux440ux435ux434ux43dux435ux435-ma1}

\[
X_n = \varepsilon_n + \theta \varepsilon_{n-1}.
\]

Ковариация: \(R(0) = \sigma^2(1+\theta^2)\),
\(R(\pm 1) = \sigma^2\theta\), \(R(|k|>1) = 0\).

Спектральная плотность:

\[
f(\lambda) = \frac{1}{2\pi} \left( \sigma^2(1+\theta^2) + 2\sigma^2\theta \cos\lambda \right) = \frac{\sigma^2}{2\pi} |1 + \theta e^{-i\lambda}|^2.
\]

\subsection{6. Контрпримеры и
предупреждения}\label{ux43aux43eux43dux442ux440ux43fux440ux438ux43cux435ux440ux44b-ux438-ux43fux440ux435ux434ux443ux43fux440ux435ux436ux434ux435ux43dux438ux44f}

\textbf{Не каждый процесс имеет плотность.} Ошибка думать, что
спектральная плотность существует всегда. Если в процессе есть строго
периодические неразрушаемые компоненты (как в Примере 2), их спектр
дискретный. Формула
\(f(\lambda) = \frac{1}{2\pi}\sum R(k) e^{-ik\lambda}\) в этом случае
просто не будет сходиться.

\textbf{Спектральная мера не является вероятностной.} В теореме
Бохнера-Хинчина (для характеристических функций) интеграл от меры равен
\(\varphi(0) = 1\), то есть мера вероятностная. В теореме Герглотца
интеграл от меры равен \(R(0) = \operatorname{Var}(X)\), что обычно не
равно 1.

\textbf{Случайная мера \(Z\) комплекснозначна даже для вещественного
процесса.} Ортогональная мера \(Z\) в спектральном разложении
\(X_n = \int e^{in\lambda} Z(d\lambda)\) принимает комплексные значения,
чтобы компенсировать экспоненты \(e^{in\lambda}\). Чтобы сам процесс
\(X_n\) был вещественным, необходимо выполнение условия симметрии:

\[
Z(\Delta) = \overline{Z(-\Delta)}.
\]

\subsection{7. Главные формулы для
доски}\label{ux433ux43bux430ux432ux43dux44bux435-ux444ux43eux440ux43cux443ux43bux44b-ux434ux43bux44f-ux434ux43eux441ux43aux438}

\begin{enumerate}
\def\labelenumi{\arabic{enumi}.}
\tightlist
\item
  \textbf{Теорема Герглотца:}
\end{enumerate}

\[
R(k) = \int_{-\pi}^{\pi} e^{i k \lambda} F(d\lambda).
\]

\begin{enumerate}
\def\labelenumi{\arabic{enumi}.}
\setcounter{enumi}{1}
\tightlist
\item
  \textbf{Прямое и обратное преобразования Фурье для плотности:}
\end{enumerate}

\[
f(\lambda) = \frac{1}{2\pi} \sum_{k=-\infty}^{\infty} R(k) e^{-i k \lambda}, \quad
R(k) = \int_{-\pi}^{\pi} e^{i k \lambda} f(\lambda) d\lambda.
\]

\begin{enumerate}
\def\labelenumi{\arabic{enumi}.}
\setcounter{enumi}{2}
\tightlist
\item
  \textbf{Спектральное представление процесса:}
\end{enumerate}

\[
X_n = \int_{-\pi}^{\pi} e^{i n \lambda} Z(d\lambda).
\]

\begin{enumerate}
\def\labelenumi{\arabic{enumi}.}
\setcounter{enumi}{3}
\tightlist
\item
  \textbf{Ортогональность приращений спектральной случайной меры:}
\end{enumerate}

\[
\mathbb E[Z(\Delta_1) \overline{Z(\Delta_2)}] = F(\Delta_1 \cap \Delta_2).
\]

\begin{enumerate}
\def\labelenumi{\arabic{enumi}.}
\setcounter{enumi}{4}
\tightlist
\item
  \textbf{Белый шум:}
\end{enumerate}

\[
f(\lambda) = \frac{\sigma^2}{2\pi}.
\]

\begin{enumerate}
\def\labelenumi{\arabic{enumi}.}
\setcounter{enumi}{5}
\tightlist
\item
  \textbf{Эргодичность по среднему через спектр:}
\end{enumerate}

\[
\mathbb E|\overline X_N|^2
=
\int_{-\pi}^{\pi}
\left|
\frac1N\sum_{n=0}^{N-1}e^{in\lambda}
\right|^2F(d\lambda),
\qquad
\overline X_N\xrightarrow{L^2}0 \Longleftrightarrow F(\{0\})=0.
\]

\subsection{8. Типичные дополнительные
вопросы}\label{ux442ux438ux43fux438ux447ux43dux44bux435-ux434ux43eux43fux43eux43bux43dux438ux442ux435ux43bux44cux43dux44bux435-ux432ux43eux43fux440ux43eux441ux44b}

\textbf{Вопрос:} Почему для дискретного времени спектр ограничен
отрезком \([-\pi, \pi]\), а для непрерывного --- всей прямой
\(\mathbb{R}\)? \textbf{Ответ:} Для дискретного времени
\(n \in \mathbb{Z}\) базисные функции \(e^{i n \lambda}\) периодичны по
\(\lambda\) с периодом \(2\pi\). То есть
\(e^{i n (\lambda + 2\pi)} = e^{i n \lambda}\). Мы физически не можем
отличить частоту \(\lambda\) от \(\lambda + 2\pi\) (эффект наложения
частот, или алиасинг). Для непрерывного времени \(t \in \mathbb{R}\)
функции \(e^{i t \lambda}\) при разных \(\lambda\) линейно независимы на
всей оси.

\textbf{Вопрос:} Как по спектральной плотности найти дисперсию процесса?
\textbf{Ответ:} Проинтегрировать ее по всем частотам:
\(\operatorname{Var}(X) = R(0) = \int_{-\pi}^{\pi} f(\lambda) d\lambda\).

\textbf{Вопрос:} Может ли спектральная плотность быть отрицательной в
каких-то точках? \textbf{Ответ:} Нет, \(f(\lambda) \ge 0\). Это прямое
следствие свойства неотрицательной определенности ковариационной
функции.

\textbf{Вопрос:} Чем отличаются спектральная мера, спектральная функция
и спектральная плотность? \textbf{Ответ:} Спектральная мера \(F\) ---
основной объект. Спектральная функция
\(\mathcal F(\lambda)=F([-\pi,\lambda])\) просто записывает эту меру как
функцию распределения по частоте. Спектральная плотность \(f\)
существует только в абсолютно непрерывном случае, когда
\(F(d\lambda)=f(\lambda)d\lambda\).

\textbf{Вопрос:} Почему спектральная мера не обязана быть вероятностной?
\textbf{Ответ:} Ее полная масса равна не \(1\), а
\(R(0)=\operatorname{Var}X_0\) для центрированного процесса. Если
хочется вероятностную меру, можно нормировать \(F/R(0)\) при \(R(0)>0\),
но в спектральной теории обычно важна именно ненормированная мощность.

\textbf{Вопрос:} Что делает линейный фильтр со спектром? \textbf{Ответ:}
Если \(Y_n=\sum_k h_kX_{n-k}\) и \(H(\lambda)=\sum_kh_ke^{-ik\lambda}\),
то спектральная мера умножается на \(|H(\lambda)|^2\). При наличии
плотности это записывается как
\(f_Y(\lambda)=|H(\lambda)|^2f_X(\lambda)\).

\textbf{Вопрос:} Связь спектра и эргодичности по среднему?
\textbf{Ответ:} В Билете 24 мы вывели, что среднее
\(\frac{1}{N}\sum_{n=0}^{N-1} X_n\) сходится в \(L_2\) к нулю (если
\(m=0\)). Через спектральную меру дисперсия среднего выражается как
интеграл от ядра Фейера. В пределе среднее сходится к величине
\(Z(\{0\})\). Поэтому эргодичность эквивалентна отсутствию атома в нуле
у спектральной меры: \(F(\{0\}) = 0\).

\subsection{9. Типичные
ошибки}\label{ux442ux438ux43fux438ux447ux43dux44bux435-ux43eux448ux438ux431ux43aux438}

\begin{itemize}
\tightlist
\item
  \textbf{Путать спектральную функцию \(\mathcal F(\lambda)\) и
  спектральную плотность \(f(\lambda)\).} Забывать, что у процесса может
  не быть плотности, если мера сингулярна или дискретна.
\item
  \textbf{Ошибки в нормировке.} Часто забывают множитель \(1/2\pi\). В
  нашем курсе он стоит перед суммой для \(f(\lambda)\), а перед
  интегралом для \(R(k)\) его нет.
\item
  \textbf{Путать \(Z(\Delta)\) и \(F(\Delta)\).} \(Z\) --- это
  \emph{случайная} мера (ее значения --- случайные величины), а \(F\)
  --- \emph{детерминированная} мера (ее значения --- действительные
  числа, дисперсии).
\item
  \textbf{Считать интеграл Риманом.} Интеграл
  \(X_n = \int e^{in\lambda} Z(d\lambda)\) --- это не интеграл
  Лебега-Стилтьеса, а стохастический интеграл, понимаемый как предел в
  \(L_2(\Omega)\).
\end{itemize}

\subsection{10. Связи с другими
билетами}\label{ux441ux432ux44fux437ux438-ux441-ux434ux440ux443ux433ux438ux43cux438-ux431ux438ux43bux435ux442ux430ux43cux438}

\begin{itemize}
\tightlist
\item
  Билет 08: Характеристические функции. Теорема Герглотца --- родная
  сестра теоремы Бохнера-Хинчина. Ковариационная функция играет роль
  характеристической функции, а спектральная мера --- роль распределения
  вероятностей.
\item
  Билет 12: Гильбертовы пространства. Изометрия между \(H_X\) и
  \(L_2(F)\) --- красивейший пример применения функционального анализа в
  теории вероятностей.
\item
  Билет 17: Стационарность в широком смысле --- основное условие для
  всей спектральной теории.
\item
  Билет 20: Дельта-функции Дирака нужны для описания дискретных спектров
  (как в Примере 2).
\item
  Билет 24: Эргодичность тесно связана со спектром (атом в нуле).
\end{itemize}

\subsection{11. Минимум на
удовлетворительно}\label{ux43cux438ux43dux438ux43cux443ux43c-ux43dux430-ux443ux434ux43eux432ux43bux435ux442ux432ux43eux440ux438ux442ux435ux43bux44cux43dux43e}

Нужно знать: 1. Что такое теорема Герглотца (написать формулу
\(R(k) = \int_{-\pi}^{\pi} e^{i k \lambda} dF(\lambda)\)). 2. Что спектр
для дискретного времени ограничен: \(\lambda \in [-\pi, \pi]\). 3. Знать
формулы, связывающие \(R(k)\) и \(f(\lambda)\) (преобразование Фурье с
множителем \(1/2\pi\)). 4. Знать, что если спектральная плотность
существует, то она неотрицательна: \(f(\lambda) \ge 0\); в общем случае
есть спектральная мера. 5. Уметь привести пример белого шума: его
ковариация равна 0 везде кроме \(k=0\), а его спектральная плотность ---
константа \(f(\lambda) = \frac{\sigma^2}{2\pi}\).

\subsection{12. Минимум на
хорошо}\label{ux43cux438ux43dux438ux43cux443ux43c-ux43dux430-ux445ux43eux440ux43eux448ux43e}

К минимуму добавляется: 1. Спектральное представление самого процесса
\(X_n = \int_{-\pi}^{\pi} e^{i n \lambda} Z(d\lambda)\). 2. Понимание
того, что такое ортогональная случайная мера \(Z(\Delta)\) и ее связь с
\(F(\Delta)\) через уравнение дисперсий:
\(\mathbb E|Z(\Delta)|^2 = F(\Delta)\). 3. Вывод спектральной плотности
для авторегрессии AR(1) или скользящего среднего MA(1).

\subsection{13. Что добавить для
отлично}\label{ux447ux442ux43e-ux434ux43eux431ux430ux432ux438ux442ux44c-ux434ux43bux44f-ux43eux442ux43bux438ux447ux43dux43e}

Для отличного ответа нужно уверенно владеть математическим аппаратом: 1.
Рассказать про изометрию гильбертовых пространств \(H_X\) и \(L_2(F)\).
2. Описать схему доказательства теоремы Герглотца (например, через
приближение функциями \(f_N(\lambda)\) и теорему Хелли/Прохорова). 3.
Привести пример с дискретным спектром (гармоническое колебание) и
объяснить, почему там нет плотности. 4. Объяснить эффект действия
линейного фильтра на спектральную плотность.

\subsection{14. Устная версия ответа на 5
минут}\label{ux443ux441ux442ux43dux430ux44f-ux432ux435ux440ux441ux438ux44f-ux43eux442ux432ux435ux442ux430-ux43dux430-5-ux43cux438ux43dux443ux442}

Спектральная теория позволяет посмотреть на стационарный в широком
смысле процесс как на сумму гармоник с разными частотами и случайными
амплитудами.

Так как время \(n\) у нас целое, все доступные частоты умещаются на
отрезке \([-\pi, \pi]\) --- вне его синусы и косинусы начнут
повторяться. Главный результат теории --- теорема Герглотца. Она
говорит, что ковариационная функция \(R(k)\) может быть представлена как
интеграл Фурье от некоторой конечной меры \(F\) на отрезке
\([-\pi, \pi]\): \emph{(пишем на доске формулу Герглотца)}. Мера \(F\)
называется спектральной мерой. Если у нее есть плотность \(f(\lambda)\),
то она неотрицательна, и \(R(k)\) с \(f(\lambda)\) связаны парой взаимно
обратных преобразований Фурье.

Кроме того, по теореме Карунена мы можем разложить в такой же интеграл и
сам процесс: \emph{(пишем интеграл
\(X_n = \int e^{in\lambda} Z(d\lambda)\))}. Здесь \(Z\) --- это уже не
обычная мера, а случайная функция множества с ортогональными
приращениями. Её дисперсия для любого интервала \(\Delta\) в точности
равна значению спектральной меры \(F(\Delta)\). Это представление
порождает идеальный словарь перевода: случайные величины из \(L_2\)
переводятся в детерминированные функции из \(L_2(F)\). Скалярные
произведения сохраняются.

Простейший пример --- белый шум. У него ковариация --- это ``иголочка''
в нуле. Значит, его спектральная плотность --- это константа. Все
частоты представлены равномерно. Если же процесс --- это просто синус со
случайной фазой, то его спектр дискретный --- это две дельта-функции на
частоте синуса.

\subsection{15. Если осталось 2 минуты перед
экзаменом}\label{ux435ux441ux43bux438-ux43eux441ux442ux430ux43bux43eux441ux44c-2-ux43cux438ux43dux443ux442ux44b-ux43fux435ux440ux435ux434-ux44dux43aux437ux430ux43cux435ux43dux43eux43c}

\begin{itemize}
\tightlist
\item
  \textbf{Время:} дискретно (\(n \in \mathbb{Z}\)). Значит, частоты
  ограничены (\(\lambda \in [-\pi, \pi]\)).
\item
  \textbf{Теорема Герглотца:}
  \(R(k) = \int_{-\pi}^{\pi} e^{i k \lambda} F(d\lambda)\). Связывает
  ковариацию и спектральную меру (дисперсию на частотах).
\item
  \textbf{Плотность:}
  \(f(\lambda) = \frac{1}{2\pi} \sum R(k) e^{-i k \lambda}\). Всегда
  \(\ge 0\).
\item
  \textbf{Сам процесс:}
  \(X_n = \int_{-\pi}^{\pi} e^{i n \lambda} Z(d\lambda)\). \(Z\) ---
  случайная ортогональная мера.
\item
  \textbf{Связь мер:} \(\mathbb E|Z(\Delta)|^2 = F(\Delta)\).
\item
  \textbf{Белый шум:}
  \(R(k) = \delta_{k0} \sigma^2 \implies f(\lambda) = \sigma^2 / 2\pi\)
  (константа).
\item
  \textbf{Ошибка:} путать \(F\) (обычная мера, задает распределение
  дисперсии) и \(Z\) (случайная мера).
\end{itemize}

\newpage

\end{document}
